\blankpage

\thispagestyle{empty}
\clearpage
\begin{tikzpicture}[remember picture, overlay]
\node at (3,-7) {\includegraphics[width=8cm]{./recettes/edités/Tarte suisse de coco (versão 1).jpg}}; 
\end{tikzpicture}\clearpage\thispagestyle{empty}
% \begin{tikzpicture}[remember picture,overlay]
% \fill[black] (-3,3.5) rectangle (14,-21.5);
% \end{tikzpicture}
\clearpage

\chapter*{Tarte suisse de coco}

\section{Ingrédients originaux}

\begin{itemize}
\item 1 verre noix de coco
\item 5 œufs
\item 2 verres farine
\item 2 verres sucre
\item 1 verre lait
\item 1 paquet beurre
\item 1 c. à s. baking
\item Vanille
\end{itemize}

\section{Préparation adaptée}

Travaillez le beurre (150--200 g.) avec 1 verre de sucre ; ajoutez les œufs, puis le lait, la vanille et la noix de coco râpée ; incorporez la farine jusqu’à obtenir une pâte homogène et ajoutez la levure. Faites cuire dans un moule beurré à 180°,\textsc{c} pendant 35--45 min. ; préparez un sirop avec 1 verre de sucre, de l’eau et du citron, portez à ébullition et arrosez le gâteau encore chaud en laissant absorber.

\chapter{Torta suíça de coco}

\section{Ingredientes originais}

\begin{itemize}
\item 1 copo de coco ralado
\item 5 ovos
\item 2 copos de farinha
\item 2 copos de açúcar
\item 1 copo de leite
\item 1 tablete de manteiga
\item 1 colher de sopa de fermento em pó
\item Baunilha
\end{itemize}

\section{Modo de preparo adaptado}

Bata a manteiga (150--200 g.) com 1 copo de açúcar; adicione os ovos, depois o leite, a baunilha e o coco ralado; incorpore a farinha até obter uma massa homogênea e misture o fermento. Asse em forma untada a 180°,\textsc{c} por 35–45 min.; prepare uma calda com 1 copo de açúcar, água e limão, ferva e regue o bolo ainda quente, deixando absorver.